\documentclass[10pt,a4paper]{amsart}
\usepackage{amsmath,amssymb,amsthm,graphicx}
\usepackage{verbatim}
\newenvironment{reference}{\comment}{\endcomment}
%\newenvironment{reference}{}{}
\newenvironment{slogan}{\comment}{\endcomment}
\newenvironment{history}{\comment}{\endcomment}

% For commutative diagrams we use Xy-pic
\usepackage[all]{xy}

% We use 2cell for 2-commutative diagrams.
\xyoption{2cell}
% Theorem environments.
%
\theoremstyle{plain}
\newtheorem{theorem}[subsection]{Theorem}
\newtheorem{proposition}[subsection]{Proposition}
\newtheorem{lemma}[subsection]{Lemma}

\theoremstyle{definition}
\newtheorem{definition}[subsection]{Definition}
\newtheorem{example}[subsection]{Example}
\newtheorem{exercise}[subsection]{Exercise}
\newtheorem{situation}[subsection]{Situation}

\theoremstyle{remark}
\newtheorem{remark}[subsection]{Remark}
\newtheorem{remarks}[subsection]{Remarks}

\numberwithin{equation}{subsection}

% Macros
%
\def\lim{\mathop{\mathrm{lim}}\nolimits}
\def\colim{\mathop{\mathrm{colim}}\nolimits}
\def\Spec{\mathop{\mathrm{Spec}}}
\def\Hom{\mathop{\mathrm{Hom}}\nolimits}
\def\Ext{\mathop{\mathrm{Ext}}\nolimits}
\def\SheafHom{\mathop{\mathcal{H}\!\mathit{om}}\nolimits}
\def\SheafExt{\mathop{\mathcal{E}\!\mathit{xt}}\nolimits}
\def\Sch{\mathit{Sch}}
\def\Mor{\mathop{\mathrm{Mor}}\nolimits}
\def\Ob{\mathop{\mathrm{Ob}}\nolimits}
\def\Sh{\mathop{\mathit{Sh}}\nolimits}
\def\NL{\mathop{N\!L}\nolimits}
\def\CH{\mathop{\mathrm{CH}}\nolimits}
\def\proetale{{pro\text{-}\acute{e}tale}}
\def\etale{{\acute{e}tale}}
\def\QCoh{\mathit{QCoh}}
\def\Ker{\mathop{\mathrm{Ker}}}
\def\Im{\mathop{\mathrm{Im}}}
\def\Coker{\mathop{\mathrm{Coker}}}
\def\Coim{\mathop{\mathrm{Coim}}}

% Boxtimes
%
\DeclareMathSymbol{\boxtimes}{\mathbin}{AMSa}{"02}

%
% Macros for moduli stacks/spaces
%
\def\QCohstack{\mathcal{QC}\!\mathit{oh}}
\def\Cohstack{\mathcal{C}\!\mathit{oh}}
\def\Spacesstack{\mathcal{S}\!\mathit{paces}}
\def\Quotfunctor{\mathrm{Quot}}
\def\Hilbfunctor{\mathrm{Hilb}}
\def\Curvesstack{\mathcal{C}\!\mathit{urves}}
\def\Polarizedstack{\mathcal{P}\!\mathit{olarized}}
\def\Complexesstack{\mathcal{C}\!\mathit{omplexes}}
% \Pic is the operator that assigns to X its picard group, usage \Pic(X)
% \Picardstack_{X/B} denotes the Picard stack of X over B
% \Picardfunctor_{X/B} denotes the Picard functor of X over B
\def\Pic{\mathop{\mathrm{Pic}}\nolimits}
\def\Picardstack{\mathcal{P}\!\mathit{ic}}
\def\Picardfunctor{\mathrm{Pic}}
\def\Deformationcategory{\mathcal{D}\!\mathit{ef}}
\usepackage{cleveref}
\crefname{originaldefinition}{original definition}{original definitions}
\Crefname{originaldefinition}{Original definition}{Original definitions}
\makeatletter
\newlabel{definition-relative-dimension}{{00QD}{source}{}{sourcepointer-1}{}}
\newlabel{definition-relative-dimension@cref}{{[originaldefinition][0][]00QD}{[1][0][]source}{}{}{}}
\makeatother
\crefname{originallemma}{original lemma}{original lemmas}
\Crefname{originallemma}{Original lemma}{Original lemmas}
\makeatletter
\newlabel{lemma-CM-over-regular-flat}{{00R4}{source}{}{sourcepointer-2}{}}
\newlabel{lemma-CM-over-regular-flat@cref}{{[originallemma][0][]00R4}{[1][0][]source}{}{}{}}
\makeatother
\crefname{originallemma}{original lemma}{original lemmas}
\Crefname{originallemma}{Original lemma}{Original lemmas}
\makeatletter
\newlabel{lemma-dimension-at-a-point-finite-type-field}{{00P1}{source}{}{sourcepointer-3}{}}
\newlabel{lemma-dimension-at-a-point-finite-type-field@cref}{{[originallemma][0][]00P1}{[1][0][]source}{}{}{}}
\makeatother
\crefname{originallemma}{original lemma}{original lemmas}
\Crefname{originallemma}{Original lemma}{Original lemmas}
\makeatletter
\newlabel{lemma-dimension-inequality-quasi-finite}{{00QF}{source}{}{sourcepointer-4}{}}
\newlabel{lemma-dimension-inequality-quasi-finite@cref}{{[originallemma][0][]00QF}{[1][0][]source}{}{}{}}
\makeatother
\crefname{originallemma}{original lemma}{original lemmas}
\Crefname{originallemma}{Original lemma}{Original lemmas}
\makeatletter
\newlabel{lemma-finite-flat-over-regular-CM}{{00R5}{source}{}{sourcepointer-5}{}}
\newlabel{lemma-finite-flat-over-regular-CM@cref}{{[originallemma][0][]00R5}{[1][0][]source}{}{}{}}
\makeatother
\title{Flatness over a quasi-finite polynomial base equals Cohen-Macaulayness of the expected local dimension}
\author{The Stacks Project Authors}
\date{stacks-project:algebra.tex:local-snapshot-20261012; GNU FDL 1.2 or later}
\begin{document}
\maketitle
\noindent\textit{Editorial note.} All delivered mathematical text is the author's own, copied from the pinned Stacks Project chapter file algebra.tex; no statement or proof step was rewritten, reordered or removed.\par\medskip
\noindent\textit{Editorial note.} Preamble provenance, as required for this source: the frozen snapshot stages the chapter's own preamble file beside it (preamble.tex, 217 lines, 8116 bytes, the same shared Stacks-wide preamble the published chow and etale-cohomology cards use); the first line of the pinned chapter file is \textbackslash{}input\{preamble\}, so no preamble text lives inside the chapter itself. The card takes that file's lines 8-18 (the verbatim package and the three \textbackslash{}comment-based environments reference/slogan/history at lines 9-12, section 13.2, together with the Xy-pic setup) and its lines 154 to the runtime end of the file (the native theorem declarations and mathematical macros). Not one character was changed, no package was added, no macro was defined, and the excerpt serves only as preamble, never as a span coordinate.\par\medskip
\noindent\textit{Editorial note.} The same-chapter pointers are located at their \textbackslash{}label lines inside the pinned chapter file; the cross-chapter pointers are located in the project's own tag list staged in the snapshot (sources/stacks-chow/tags.txt), at the verbatim lines carrying their canonical tags.\par\medskip
\noindent\textit{Editorial note.} The counted statement refers to the definition of relative dimension at a point; it is recorded as a navigation stub with kind definition, and the argument using it stays inside the counted span.\par\medskip
\noindent\textit{Editorial note.} The counted proof also refers to same-chapter statements by label, recorded the same way; no \textbackslash{}cite appears, so the card carries no bibliography.\par\medskip
\noindent\textit{Allowed standard prerequisites (editorial).}\begin{itemize}
\item The dimension inequality for a quasi-finite map. Invoked in the counted proof at lines 33759-33761. Bounds the dimension of the local ring by that of the polynomial base.
\item Flatness of a Cohen-Macaulay module over a regular local ring. Invoked at the end of the counted proof (lines 33783-33785). Converts the Cohen-Macaulayness and dimension condition back into flatness.
\end{itemize}\medskip
\ifdefined\hypertarget\hypertarget{sourcepointer-1}{}\fi
\noindent\textit{Source reference.} Original definition 00QD: sources/stacks-algebra/pinned-original-source.tex, source lines 31367--31367. Navigation pointer for \textbackslash{}ref\{definition-relative-dimension\} used inside the counted span (line 33753); the argument that consumes it is retained in the span.\par\medskip
\ifdefined\hypertarget\hypertarget{sourcepointer-2}{}\fi
\noindent\textit{Source reference.} Original lemma 00R4: sources/stacks-algebra/pinned-original-source.tex, source lines 32871--32871. Navigation pointer for \textbackslash{}ref\{lemma-CM-over-regular-flat\} used inside the counted span (line 33783); the argument that consumes it is retained in the span.\par\medskip
\ifdefined\hypertarget\hypertarget{sourcepointer-3}{}\fi
\noindent\textit{Source reference.} Original lemma 00P1: sources/stacks-algebra/pinned-original-source.tex, source lines 28439--28439. Navigation pointer for \textbackslash{}ref\{lemma-dimension-at-a-point-finite-type-field\} used inside the counted span (line 33772); the argument that consumes it is retained in the span.\par\medskip
\ifdefined\hypertarget\hypertarget{sourcepointer-4}{}\fi
\noindent\textit{Source reference.} Original lemma 00QF: sources/stacks-algebra/pinned-original-source.tex, source lines 31484--31484. Navigation pointer for \textbackslash{}ref\{lemma-dimension-inequality-quasi-finite\} used inside the counted span (line 33761); the argument that consumes it is retained in the span.\par\medskip
\ifdefined\hypertarget\hypertarget{sourcepointer-5}{}\fi
\noindent\textit{Source reference.} Original lemma 00R5: sources/stacks-algebra/pinned-original-source.tex, source lines 27706--27706. Navigation pointer for \textbackslash{}ref\{lemma-finite-flat-over-regular-CM\} used inside the counted span (line 33768); the argument that consumes it is retained in the span.\par\medskip
\begin{lemma}
\label{lemma-where-CM}
Let $S$ be a finite type algebra over a field $k$.
Let $\varphi : k[y_1, \ldots, y_d] \to S$ be a quasi-finite ring map.
As subsets of $\Spec(S)$ we have
$$
\{ \mathfrak q \mid
S_{\mathfrak q} \text{ flat over }k[y_1, \ldots, y_d]\}
=
\{ \mathfrak q \mid
S_{\mathfrak q} \text{ CM and }\dim_{\mathfrak q}(S/k) = d\}
$$
For notation see Definition \ref{definition-relative-dimension}.
\end{lemma}

\begin{proof}
Let $\mathfrak q \subset S$ be a prime. Denote
$\mathfrak p = k[y_1, \ldots, y_d] \cap \mathfrak q$.
Note that always
$\dim(S_{\mathfrak q}) \leq \dim(k[y_1, \ldots, y_d]_{\mathfrak p})$
by Lemma \ref{lemma-dimension-inequality-quasi-finite} for example.
Moreover, the field extension $\kappa(\mathfrak q)/\kappa(\mathfrak p)$
is finite and hence
$\text{trdeg}_k(\kappa(\mathfrak p)) = \text{trdeg}_k(\kappa(\mathfrak q))$.

\medskip\noindent
Let $\mathfrak q$ be an element of the left hand side.
Then Lemma \ref{lemma-finite-flat-over-regular-CM} applies
and we conclude that $S_{\mathfrak q}$ is Cohen-Macaulay
and $\dim(S_{\mathfrak q}) = \dim(k[y_1, \ldots, y_d]_{\mathfrak p})$.
Combined with the equality of transcendence degrees above and
Lemma \ref{lemma-dimension-at-a-point-finite-type-field} this
implies that $\dim_{\mathfrak q}(S/k) = d$. Hence $\mathfrak q$
is an element of the right hand side.

\medskip\noindent
Let $\mathfrak q$ be an element of the right hand side.
By the equality of transcendence degrees above, the assumption
that $\dim_{\mathfrak q}(S/k) = d$ and
Lemma \ref{lemma-dimension-at-a-point-finite-type-field}
we conclude that
$\dim(S_{\mathfrak q}) = \dim(k[y_1, \ldots, y_d]_{\mathfrak p})$.
Hence Lemma \ref{lemma-CM-over-regular-flat}
applies and we see that $\mathfrak q$ is an
element of the left hand side.
\end{proof}

\end{document}
